% ============================================================
% Appendix B: Glossary of Key Terms
% ============================================================

\chapter*{Glossary}
\addcontentsline{toc}{chapter}{Glossary of Key Terms}

\begin{description}
    \item[Algebraic Number] A number that is a root of a nonzero polynomial with rational (or equivalently, integer) coefficients. For example, $\sqrt{2}$ is algebraic because it is a root of $x^2 - 2 = 0$. Numbers that are not algebraic are called \textit{transcendental} (e.g., $\pi$ and $e$).

    \item[Axiom System] A set of foundational assumptions (axioms) from which all theorems of a theory are derived. An axiom system is \textit{consistent} if no contradiction can be derived, and \textit{complete} if every statement can be decided.

    \item[Betti Numbers] Topological invariants $b_0, b_1, b_2, \ldots$ of a space, where $b_k$ counts the number of $k$-dimensional ``holes.'' For a surface of genus $g$: $b_0 = 1$, $b_1 = 2g$, $b_2 = 1$.

    \item[Cardinality] A measure of the ``size'' of a set. Two sets have the same cardinality if there is a bijection between them. The cardinality of the natural numbers is $\aleph_0$; the cardinality of the continuum is $\mathfrak{c} = 2^{\aleph_0}$.

    \item[Class Field Theory] The theory describing abelian extensions of number fields in terms of their arithmetic properties. Central results include the Artin reciprocity theorem and the Kronecker--Weber theorem.

    \item[Continuum Hypothesis (CH)] The conjecture that there is no set whose cardinality is strictly between that of the natural numbers and the real numbers. Equivalently: $\mathfrak{c} = \aleph_1$. Proved to be independent of ZFC.

    \item[Convergence] A sequence $(a_n)$ converges to $L$ if for every $\epsilon > 0$, there exists $N$ such that $|a_n - L| < \epsilon$ for all $n > N$.

    \item[Covering Space] A space $\tilde{X}$ together with a continuous surjection $p: \tilde{X} \to X$ such that every point of $X$ has an open neighborhood evenly covered by $p$.

    \item[Dehn Invariant] An invariant of polyhedra under scissors congruence. For a polyhedron $P$, it is $D(P) = \sum_i \ell_i \otimes \theta_i \in \mathbb{R} \otimes (\mathbb{R}/\pi\mathbb{Z})$, where the sum is over all edges with length $\ell_i$ and dihedral angle $\theta_i$.

    \item[Diophantine Equation] A polynomial equation for which we seek integer (or rational) solutions. For example, $x^2 + y^2 = z^2$ (Pythagorean triples) and $x^n + y^n = z^n$ (Fermat's Last Theorem).

    \item[Dirichlet Problem] Given a bounded domain $\Omega \subset \mathbb{R}^n$ and a function $f$ on $\partial\Omega$, find a function $u$ that is harmonic in $\Omega$ and equals $f$ on $\partial\Omega$.

    \item[Elliptic PDE] A partial differential equation whose principal part is positive definite. The canonical example is Laplace's equation: $\Delta u = 0$. Solutions to elliptic PDEs are typically smooth.

    \item[Formal System] A system of axioms and inference rules. A formal system is \textit{consistent} if no contradiction can be derived, and \textit{complete} if every statement is either provable or disprovable.

    \item[Forcing] A technique developed by Paul Cohen (1963) for constructing models of set theory in which certain statements hold or fail. Cohen used forcing to prove the independence of the Continuum Hypothesis.

    \item[Fuchsian Group] A discrete subgroup of $\operatorname{PSL}(2, \mathbb{R})$, the group of orientation-preserving isometries of the hyperbolic plane. Fuchsian groups arise in the uniformization of Riemann surfaces.

    \item[Galois Group] For a field extension $L/K$, the group of field automorphisms of $L$ that fix $K$ pointwise. The fundamental theorem of Galois theory connects subgroups of the Galois group to intermediate fields.

    \item[Geodesic] The ``shortest path'' between two points in a curved space. In Euclidean space, geodesics are straight lines; on a sphere, they are great circles.

    \item[Goldbach Conjecture] Every even integer greater than 2 can be expressed as the sum of two primes. Proposed by Christian Goldbach in 1742. Still unsolved.

    \item[Harmonic Function] A function $u: \Omega \to \mathbb{R}$ satisfying Laplace's equation $\Delta u = 0$, where $\Delta$ is the Laplacian operator.

    \item[Hilbert Space] A complete inner product space. The key infinite-dimensional analog of Euclidean space, central to quantum mechanics and functional analysis.

    \item[Incompleteness Theorem (G\"odel)] (1) Any consistent, recursively axiomatized system containing basic arithmetic is incomplete---there exist statements that are neither provable nor disprovable. (2) Such a system cannot prove its own consistency.

    \item[Lie Group] A group that is also a smooth manifold, where the group operations are smooth. Examples: $\operatorname{GL}(n, \mathbb{R})$, $\operatorname{SO}(n)$, the circle group $S^1$.

    \item[Lie Algebra] The tangent space at the identity of a Lie group, equipped with the Lie bracket $[X, Y] = XY - YX$. The Lie algebra encodes the local structure of the Lie group.

    \item[Local-Global Principle] The principle that a global property (e.g., a quadratic form representing 0 over $\mathbb{Q}$) can be decided by checking it locally (at each prime $p$ and at infinity). The Hasse--Minkowski theorem is a precise instance.

    \item[Manifold] A topological space that locally resembles Euclidean space. Formally, a space where every point has a neighborhood homeomorphic to an open subset of $\mathbb{R}^n$.

    \item[Metric Space] A set $X$ equipped with a distance function $d: X \times X \to \mathbb{R}$ satisfying positivity, symmetry, and the triangle inequality.

    \item[Monodromy] The behavior of a multivalued function (e.g., $\log z$, $\sqrt{z}$) when analytically continued along a closed loop. The monodromy theorem states that in a simply connected domain, continuation along any closed loop returns to the original value.

    \item[Poincar\'e Conjecture] Every simply connected, closed 3-manifold is homeomorphic to the 3-sphere $S^3$. Proved by Grigori Perelman in 2003.

    \item[Quadratic Form] A homogeneous polynomial of degree 2, e.g., $Q(x, y) = ax^2 + bxy + cy^2$. The theory of quadratic forms studies how these polynomials represent numbers.

    \item[Riemann Surface] A one-dimensional complex manifold---a surface with a complex-analytic structure. Examples: the Riemann sphere $\widehat{\mathbb{C}}$, the complex plane $\mathbb{C}$, a torus $\mathbb{C}/\Lambda$.

    \item[Scissors Congruence] Two polyhedra are scissors congruent if one can be cut into finitely many pieces and reassembled to form the other. Volume is a complete invariant in 2D; in 3D, the Dehn invariant provides an additional obstruction.

    \item[Sobolev Space] A function space equipped with a norm involving both the function and its weak derivatives. Central to the modern theory of PDEs.

    \item[Tessellation] A covering of a space by shapes (tiles) with no overlaps. Regular tessellations are tilings by regular polygons (in 2D) or regular polytopes (in higher dimensions).

    \item[Transcendental Number] A number that is not algebraic---i.e., not the root of any nonzero polynomial with rational coefficients. Examples: $\pi$, $e$, $e^{\pi}$.

    \item[Uniformization Theorem] Every simply connected Riemann surface is conformally equivalent to one of three surfaces: the Riemann sphere $\widehat{\mathbb{C}}$, the complex plane $\mathbb{C}$, or the unit disk $\mathbb{D}$.

    \item[ZFC] Zermelo--Fraenkel set theory with the Axiom of Choice. The standard axiom system for mathematics. Hilbert's second problem asked about the consistency of arithmetic within such a system.
\end{description}
